\section{Sections of sheaves of modules}
\label{section-sections}

\noindent
Let $(X, \mathcal{O}_X)$ be a ringed space.
Let $\mathcal{F}$ be a sheaf of $\mathcal{O}_X$-modules.
Let $s \in \Gamma(X, \mathcal{F}) = \mathcal{F}(X)$ be a
global section. There is a unique {\it map of $\mathcal{O}_X$-modules}
$$
\mathcal{O}_X \longrightarrow \mathcal{F}, \ f \longmapsto fs
$$
{\it associated to $s$}. The notation above signifies that a local
section $f$ of $\mathcal{O}_X$, i.e., a section $f$ over some open $U$,
is mapped to the multiplication of $f$ with the restriction of $s$ to
$U$. Conversely, any map $\varphi : \mathcal{O}_X \to \mathcal{F}$
gives rise to a section $s = \varphi(1)$ such that $\varphi$ is
the morphism associated to $s$.

\begin{definition}
\label{definition-globally-generated}
Let $(X, \mathcal{O}_X)$ be a ringed space.
Let $\mathcal{F}$ be a sheaf of $\mathcal{O}_X$-modules.
We say that $\mathcal{F}$ is {\it generated by global
sections} if there exist a set $I$, and
global sections $s_i \in \Gamma(X, \mathcal{F})$, $i \in I$
such that the map
$$
\bigoplus\nolimits_{i \in I}
\mathcal{O}_X \longrightarrow \mathcal{F}
$$
which is the map associated to $s_i$ on the summand corresponding to $i$,
is surjective. In this case we say that the sections $s_i$
{\it generate} $\mathcal{F}$.
\end{definition}

\noindent
We often use the abuse of notation introduced in
Sheaves, Section \ref{sheaves-section-stalks} where, given a local
section $s$ of $\mathcal{F}$ defined in an open neighbourhood
of a point $x \in X$, we denote $s_x$, or even $s$ the image of $s$
in the stalk $\mathcal{F}_x$.

\begin{lemma}
\label{lemma-globally-generated}
Let $(X, \mathcal{O}_X)$ be a ringed space.
Let $\mathcal{F}$ be a sheaf of $\mathcal{O}_X$-modules.
Let $I$ be a set. Let
$s_i \in \Gamma(X, \mathcal{F})$, $i \in I$
be global sections. The sections $s_i$ generate
$\mathcal{F}$ if and only if for all $x\in X$ the
elements $s_{i, x} \in \mathcal{F}_x$ generate
the $\mathcal{O}_{X, x}$-module $\mathcal{F}_x$.
\end{lemma}

\begin{proof}
Omitted.
\end{proof}

\begin{lemma}
\label{lemma-tensor-product-globally-generated}
\begin{slogan}
The tensor product of globally generated sheaves of modules is
globally generated.
\end{slogan}
Let $(X, \mathcal{O}_X)$ be a ringed space.
Let $\mathcal{F}$, $\mathcal{G}$ be sheaves of $\mathcal{O}_X$-modules.
If $\mathcal{F}$ and $\mathcal{G}$ are generated by global sections
then so is $\mathcal{F} \otimes_{\mathcal{O}_X} \mathcal{G}$.
\end{lemma}

\begin{proof}
Omitted.
\end{proof}

\begin{lemma}
\label{lemma-generated-by-local-sections}
Let $(X, \mathcal{O}_X)$ be a ringed space.
Let $\mathcal{F}$ be a sheaf of $\mathcal{O}_X$-modules.
Let $I$ be a set. Let $s_i$, $i \in I$ be a collection
of local sections of $\mathcal{F}$, i.e., $s_i \in \mathcal{F}(U_i)$
for some opens $U_i \subset X$. There exists a unique smallest
subsheaf of $\mathcal{O}_X$-modules $\mathcal{G}$ such
that each $s_i$ corresponds to a local section of
$\mathcal{G}$.
\end{lemma}

\begin{proof}
Consider the subpresheaf of $\mathcal{O}_X$-modules
defined by the rule
$$
U
\longmapsto
\{
\text{sums } \sum\nolimits_{i \in J} f_i (s_i|_U)
\text{ where } J \text{ is finite, }
U \subset U_i \text{ for } i\in J, \text{ and }
f_i \in \mathcal{O}_X(U)
\}
$$
Let $\mathcal{G}$ be the sheafification of this subpresheaf.
This is a subsheaf of $\mathcal{F}$ by
Sheaves, Lemma \ref{sheaves-lemma-sheafify-epi-mono}.
Since all the finite sums clearly have to be in $\mathcal{G}$
this is the smallest subsheaf as desired.
\end{proof}

\begin{definition}
\label{definition-generated-by-local-sections}
Let $(X, \mathcal{O}_X)$ be a ringed space.
Let $\mathcal{F}$ be a sheaf of $\mathcal{O}_X$-modules.
Given a set $I$, and
local sections $s_i$, $i \in I$ of $\mathcal{F}$
we say that the subsheaf $\mathcal{G}$ of
Lemma \ref{lemma-generated-by-local-sections}
above is the {\it subsheaf generated by the $s_i$}.
\end{definition}

\begin{lemma}
\label{lemma-generated-by-local-sections-stalk}
Let $(X, \mathcal{O}_X)$ be a ringed space.
Let $\mathcal{F}$ be a sheaf of $\mathcal{O}_X$-modules.
Given a set $I$, and
local sections $s_i$, $i \in I$ of $\mathcal{F}$.
Let $\mathcal{G}$ be the subsheaf generated by the
$s_i$ and let $x\in X$.
Then $\mathcal{G}_x$ is the $\mathcal{O}_{X, x}$-submodule of
$\mathcal{F}_x$ generated by the elements $s_{i, x}$
for those $i$ such that $s_i$ is defined at $x$.
\end{lemma}

\begin{proof}
This is clear from the construction of $\mathcal{G}$
in the proof of Lemma \ref{lemma-generated-by-local-sections}.
\end{proof}










